% Examples shared by the Z statements and their individual proofs.
\newcommand{\Zexampleone}{%
  \textbf{Z1: at most three switches}\par\medskip
  \begin{shiftchain}[
    column sep=4.5mm, row sep=6mm, dot radius=3pt,
    labels=false, label font=\scriptsize,
    switches={3,5,6},
    switch labels={3/$y_{i-1}$,5/$y_i$,6/\raisebox{-3mm}{$z$}},
  ]{9}
    \shiftedge{1/cfred,2/cfred,6/cfred,9/cfnavy}[$\vec e$]
  \end{shiftchain}
  \par\medskip Two old switches, one new switch.
}

\newcommand{\Zexampletwo}{%
  \textbf{Z2: repair the chosen edge}\par\medskip
  \begin{shiftchain}[
    column sep=4.5mm, row sep=6mm, dot radius=3pt,
    labels=false, label font=\scriptsize,
    switches={3,5,6/dashed},
    switch labels={3/$y_{i-1}$,5/$y_i$,6/\raisebox{-3mm}{$\hat z$}},
  ]{9}
    \shiftedge{1/cfred,2/cfred,6/cfred,9/cfred}[$t$]
  \end{shiftchain}
  \par\medskip
  \begin{shiftchain}[
    column sep=4.5mm, row sep=6mm, dot radius=3pt,
    labels=false, label font=\scriptsize,
    switches={3,5,6},
    switch labels={3/$y_{i-1}$,5/$y_i$,6/\raisebox{-3mm}{$\hat z$}},
  ]{9}
    \shiftedge{1/cfred,2/cfred,6/cfred,9/cfnavy}[$t+1$]
  \end{shiftchain}
  \par\medskip No new monochromatic edges.
}

\newcommand{\Zexamplethree}{%
  \textbf{Z3: surviving monochromatic edges}\par\medskip
  \begin{shiftchain}[
    column sep=4mm, row sep=6mm, dot radius=3pt,
    labels=false, label font=\scriptsize,
    switches={1,5,7},
    switch labels={1/$z$,5/$y_{j-1}$,7/$y_j$},
  ]{10}
    \shiftedge{2/cfred,4/cfred,8/cfred,10/cfred}[$\vec f$]
  \end{shiftchain}
  \par\medskip $z<f_1$ for every $z\in Z(t)$.
}

\newcommand{\Zexamplefour}{%
  \textbf{Z4: two switches bracket a vertex}\par\medskip
  \begin{shiftchain}[
    column sep=5mm, row sep=6mm, dot radius=3pt,
    labels=false, label font=\scriptsize,
    switches={3,4,7/dashed},
    switch labels={3/$y_j$,4/\raisebox{-3mm}{$z$},7/$x_j$},
  ]{8}
    \node at (4,0) {$v$};
    \shiftedge{1/cfred,2/cfred,4/cfnavy,8/cfred}[$\vec f$]
  \end{shiftchain}
  \par\medskip An inserted $z$ occurs only immediately right of $v$:\par
  $y_j=v^-,\quad z=v^+\leq x_j$.
}

\begin{frame}[t]{Stage Z: recap}
  \small\setstretch{1}
  \claimreminder{\Ysix}
  Sweep left to right; insert a switch just after the third vertex
  of each remaining $(2,3)$-edge.

  \smallskip
  \centering
  \only<1>{%
  \outlinediagram
    {3,5,9,13,15,19}
    {3/$x_1$,5/$y_2$,9/$x_3$,13/$x_4$,15/$y_5$,19/$x_6$}
    {1,2,3,6,7,8,9,14,15,20}
    {4,5,10,11,12,13,16,17,18,19}
    {2}
    {Z: left to right (LR) --- $e_{2}$ is monochromatic}
    {1}%
}%
  \only<2>{%
  \outlinediagram
    {3,5,6/outline moved,9,13,15,19}
    {3/$x_1$,5/$y_2$,6/$z_1$,9/$x_3$,13/$x_4$,15/$y_5$,19/$x_6$}
    {1,2,3,6,10,11,12,13,16,17,18,19}
    {4,5,7,8,9,14,15,20}
    {2}
    {Z: left to right (LR) --- insert $z_1$ after vertex $6$}
    {1}%
}%
  \only<3>{%
  \outlinediagram
    {3,5,6,9,13,15,19}
    {3/$x_1$,5/$y_2$,6/$z_1$,9/$x_3$,13/$x_4$,15/$y_5$,19/$x_6$}
    {1,2,3,6,10,11,12,13,16,17,18,19}
    {4,5,7,8,9,14,15,20}
    {7}
    {Z: left to right (LR) --- $e_{7}$ is monochromatic}
    {1}%
}%
  \only<4>{%
  \outlinediagram
    {3,5,6,9,13,15,16/outline moved,19}
    {3/$x_1$,5/$y_2$,6/$z_1$,9/$x_3$,13/$x_4$,15/$y_5$,16/$z_2$,19/$x_6$}
    {1,2,3,6,10,11,12,13,16,20}
    {4,5,7,8,9,14,15,17,18,19}
    {7}
    {Z: left to right (LR) --- insert $z_2$ after vertex $16$}
    {1}%
}%
  \only<5>{%
  \outlinediagram
    {3,5,6,9,13,15,16,19}
    {3/$x_1$,5/$y_2$,6/$z_1$,9/$x_3$,13/$x_4$,15/$y_5$,16/$z_2$,19/$x_6$}
    {1,2,3,6,10,11,12,13,16,20}
    {4,5,7,8,9,14,15,17,18,19}
    {0}
    {Proper 2-colouring: two Y moves, two Z insertions}
    {1}%
}%
\end{frame}

% Introduce all four invariants before the induction, as in Stage Y.
\begin{frame}[t]{Z1--Z4: what Stage Z preserves and achieves}
  {\normalfont\small\color{black}\setstretch{1}
    $t$: number of Stage Z insertions; $Z(t)$: inserted switches\par
    $c_{Z,t}$: colouring with switches $Y\cup Z(t)$; $c_{Z,0}=c_Y$\par
    $M(c)$: monochromatic edges in $c$\par
    Induction on $t$: Z1--Z4; the $Y$-switches stay fixed\par
  }
  \smallskip
  \begin{columns}[T,onlytextwidth]
    \begin{column}{0.58\textwidth}
      \footnotesize\setstretch{1}
      \begin{itemize}[itemsep=7pt]
        \item<1-> \Zone
        \item<2-> \Ztwo
        \item<3-> \Zthree
        \item<4-> \Zfour
      \end{itemize}
    \end{column}
    \begin{column}{0.39\textwidth}
      \centering\small\setstretch{1}
      \begin{overlayarea}{\linewidth}{4.5cm}
        \raggedright
        \only<1>{\Zexampleone}
        \only<2>{\Zexampletwo}
        \only<3>{\Zexamplethree}
        \only<4>{\Zexamplefour}
      \end{overlayarea}
    \end{column}
  \end{columns}
\end{frame}

% All inductive hypotheses refer to time t; prove the claims at t+1 below.
\begin{frame}[t]{Z1--Z4: one Stage Z step}
  \footnotesize\setstretch{1}
  \claimreminder{\Zthree\Yfive}
  Assume Z1--Z4 at time $t$; choose a minimal monochromatic edge $\vec e$.
  Insert $\hat z=e_3^+$.
  \begin{center}
    \only<1>{%
      \begin{shiftchain}[
        column sep=6mm, row sep=6mm, dot radius=3pt,
        labels=false, label font=\scriptsize,
        switches={0,5,7,12},
        switch labels={0/$y_{i-2}$,5/$y_{i-1}$,7/$y_i$,12/$y_{i+1}$},
      ]{16}
        \node at (8,0) {$e_3$};
        \shiftedge{2/cfred,4/cfred,8/cfred,11/cfred}[$\vec e$]
        \shiftchainbounds{-2}{-3.7}{16.7}{0.9}
      \end{shiftchain}
    }
    \only<2->{%
      \begin{shiftchain}[
        column sep=6mm, row sep=6mm, dot radius=3pt,
        labels=false, label font=\scriptsize,
        switches={0,5,7,8,12},
        switch labels={0/$y_{i-2}$,5/$y_{i-1}$,7/$y_i$,8/\raisebox{-3mm}{$\hat z$},12/$y_{i+1}$},
      ]{16}
        \node at (8,0) {$e_3=g_3$};
        \shiftedge{2/cfred,4/cfred,8/cfred,11/cfnavy}[$\vec e$]
        \uncover<3->{\shiftedge{6/cfnavy,7/cfnavy,8/cfred,16/cfred}[$\vec g$]}
        \shiftchainbounds{-2}{-3.7}{16.7}{0.9}
      \end{shiftchain}
    }
  \end{center}
  \begin{itemize}[itemsep=4pt]
    \item<1-> Z3: $\vec e$ has two $Y$-switches in $(e_2,e_3)$;
      every old $z\in Z(t)$ lies before $e_1$.
    \item<2-> Y5: $y_i=e_3^-<x_i$; inserting $\hat z=e_3^+$ changes the colour of $e_4$.
    \item<3-> Y5 also gives $g_3=e_3$ and $y_{i+1}<g_4$.
  \end{itemize}
\end{frame}

\begin{frame}[t]{Z1: older inserted switches lie before $x_{i-2}$}
  \footnotesize\setstretch{1}
  \tcbset{before skip=4pt,after skip=4pt}
  \claimreminder{\Zthree\Zfour\Ytwo}
  \begin{lemma}{Claim}
    Every $z\in Z(t)$ satisfies $z\leq x_{i-2}$.
  \end{lemma}
  \begin{proof}{Proof}
    \begin{itemize}[itemsep=4pt]
      \item<1-> Z3 gives $z<e_1<y_{i-1}$.
      \item<2-> By Z4, $y_j=v^-<z=v^+\leq x_j$ for some $j$.
      \item<3-> Thus $j\leq i-2$, so $z\leq x_j\leq x_{i-2}$.
        Y2 also gives $x_{i-2}<y_{i-1}$.
    \end{itemize}
  \end{proof}
  \begin{center}
    \begin{shiftchain}[
      column sep=6mm, row sep=6mm, dot radius=3pt,
      labels=false, label font=\scriptsize,
      switches={2,3,5/dashed,8,11},
      switch labels={2/$y_j$,3/\raisebox{-3mm}{$z$},5/$x_{i-2}$,8/$y_{i-1}$,11/$y_i$},
    ]{16}
      \uncover<2->{\node at (3,0) {$v$};}
      \shiftedge{6/cfred,7/cfred,12/cfred,16/cfred}[$\vec e$]
      \shiftchainbounds{-2}{-2.7}{16.7}{0.9}
    \end{shiftchain}
  \end{center}
\end{frame}

\begin{frame}[t]{Z1: every edge keeps at most three switches}
  \footnotesize\setstretch{1}
  \tcbset{before skip=4pt,after skip=4pt}
  \claimreminder{\Yone\Ytwo\Xtwo\Zone
    \textbf{Previous step.} Every $z\in Z(t)$ satisfies $z\leq x_{i-2}$.}
  \begin{lemma}{Claim}\Zone\end{lemma}
  \begin{proof}{Proof}
    \begin{itemize}[itemsep=4pt]
      \item<1-> \textbf{Base.} Y1: at most two switches in $c_{Z,0}=c_Y$.
      \item<2-> \textbf{Step.} A new fourth switch in $\vec f$ requires both
        $\hat z$ and an old $z\in Z(t)$: Y1 allows only two $Y$-switches.
      \item<3-> By the previous step and Y2,
        $f_1<z\leq x_{i-2}<x_{i-1}<y_i<\hat z<f_4$.
      \item<4-> X2 forces $x_{i-1}=f_4^-$, but $x_{i-1}<\hat z<f_4$.
    \end{itemize}
  \end{proof}
  \begin{center}
    \only<2>{%
      \begin{shiftchain}[
        column sep=6mm, row sep=5mm, dot radius=3pt,
        labels=false, label font=\scriptsize,
        switches={3,8,11,12},
        switch labels={3/$z$,8/$y_{i-1}$,11/$y_i$,12/\raisebox{-3mm}{$\hat z$}},
      ]{16}
        \shiftedge{3/cfred,6/cfnavy,10/cfred,16/cfred}[$\vec f$]
        \shiftchainbounds{-2}{-2.7}{16.7}{0.9}
      \end{shiftchain}
    }
    \only<3->{%
      \begin{shiftchain}[
        column sep=6mm, row sep=5mm, dot radius=3pt,
        labels=false, label font=\scriptsize,
        switches={3,5/dashed,8,9/dashed,11,12,15/densely dotted},
        switch labels={3/$z$,5/\raisebox{-3mm}{$x_{i-2}$},8/$y_{i-1}$,9/\raisebox{-3mm}{$x_{i-1}$},11/$y_i$,12/\raisebox{-3mm}{$\hat z$},15/$f_4^-$},
      ]{16}
        \shiftedge{3/cfred,6/cfnavy,10/cfred,16/cfred}[$\vec f$]
        \shiftchainbounds{-2}{-2.7}{16.7}{0.9}
      \end{shiftchain}
    }
  \end{center}
\end{frame}

\begin{frame}[t]{Z2: a new monochromatic edge has only two possible switch pairs}
  \footnotesize\setstretch{1}
  \tcbset{before skip=4pt,after skip=4pt}
  \claimreminder{\Yone\Zone\Zthree\switchparity}
  \begin{lemma}{Claim}\Ztwo\end{lemma}
  \begin{proof}{Proof}
    \begin{itemize}[itemsep=4pt]
      \item<1-> \textbf{Base.} No insertion at $t=0$; Z2 starts at $t=1$.
      \item<2-> \textbf{Step.} The chosen $\vec e$ is no longer monochromatic.
        Suppose a new monochromatic edge $\vec f$ appears.
      \item<3-> It contains $\hat z$; otherwise all its vertices change colour together, or none do.
      \item<4-> Y1, Z1 and Parity: exactly two switches, in one gap.
      \item<5-> Z3 puts all old $z$ before $e_1$.
        The partner of $\hat z$ is therefore $y_i$ or $y_{i+1}$.
    \end{itemize}
  \end{proof}
  \begin{center}
    \uncover<2->{%
      \begin{shiftchain}[
        column sep=6mm, row sep=5mm, dot radius=3pt,
        labels=false, label font=\scriptsize,
        switches={4,7,8,12},
        switch labels={4/$y_{i-1}$,7/$y_i$,8/\raisebox{-3mm}{$\hat z$},12/$y_{i+1}$},
      ]{14}
        \node at (8,0) {$e_3$};
        \shiftedge{1/cfred,3/cfred,8/cfred,11/cfnavy}[$\vec e$]
      \end{shiftchain}
    }
  \end{center}
\end{frame}

\begin{frame}[t]{Z2: $y_i,\hat z$ cannot lie in the last gap}
  \footnotesize\setstretch{1}
  \tcbset{before skip=4pt,after skip=4pt}
  \claimreminder{\Ztwo\switchparity}
  \begin{lemma}{Claim}
    No new monochromatic edge has only $y_i,\hat z$ in $(f_3,f_4)$.
  \end{lemma}
  \begin{proof}{Proof}
    \begin{itemize}[itemsep=4pt]
      \item<1-> Suppose $f_3<y_i<\hat z<f_4$.
      \item<2-> Since $f_3<e_3$, comparability gives $\vec f<\vec e$ and $f_1\leq e_1$.
      \item<3-> But $f_1\leq e_1<y_{i-1}<y_i<\hat z<f_4$:
        a third switch lies in $\vec f$.
    \end{itemize}
  \end{proof}
  \begin{center}
    \only<1>{%
      \begin{shiftchain}[
        column sep=6mm, row sep=6mm, dot radius=3pt,
        labels=false, label font=\scriptsize,
        switches={9,10}, switch labels={9/$y_i$,10/\raisebox{-3mm}{$\hat z$}},
      ]{16}
        \node at (5,0) {$f_3$};
        \shiftedge{1/cfred,3/cfred,5/cfred,12/cfred}[$\vec f$]
        \shiftchainbounds{-2}{-3.7}{16.7}{0.9}
      \end{shiftchain}
    }
    \only<2->{%
      \begin{shiftchain}[
        column sep=6mm, row sep=6mm, dot radius=3pt,
        labels=false, label font=\scriptsize,
        switches={7,9,10},
        switch labels={7/$y_{i-1}$,9/$y_i$,10/\raisebox{-3mm}{$\hat z$}},
      ]{16}
        \node at (5,0) {$f_3$};
        \node at (10,0) {$e_3$};
        \shiftedge{1/cfred,3/cfred,5/cfred,12/cfnavy}[$\vec f$]
        \shiftedge{4/cfred,6/cfred,10/cfred,14/cfnavy}[$\vec e$]
        \shiftchainbounds{-2}{-3.7}{16.7}{0.9}
      \end{shiftchain}
    }
  \end{center}
\end{frame}

\begin{frame}[t]{Z2: $y_i,\hat z$ cannot lie in either earlier gap}
  \footnotesize\setstretch{1}
  \tcbset{before skip=4pt,after skip=4pt}
  \claimreminder{\Yfive\Ztwo\switchparity}
  \begin{lemma}{Claim}
    No new monochromatic edge has only $y_i,\hat z$ in $(f_1,f_2)$ or $(f_2,f_3)$.
  \end{lemma}
  \begin{proof}{Proof}
    \begin{itemize}[itemsep=4pt]
      \item<1-> In both cases, $\hat z<f_3$.
      \item<2-> Y5 gives $g_3=e_3<\hat z<f_3$ and $y_{i+1}<g_4$.
      \item<3-> Comparability gives $\vec g<\vec f$, so $y_{i+1}<g_4\leq f_4$:
        a third switch lies in $\vec f$.
    \end{itemize}
  \end{proof}
  \begin{columns}[T,onlytextwidth]
    \begin{column}{0.48\textwidth}
      \centering\scriptsize
      \textbf{Gap $(f_1,f_2)$}\par\smallskip
      \only<1>{%
        \begin{shiftchain}[
          column sep=3.5mm, row sep=5mm, dot radius=2.7pt,
          labels=false, label font=\tiny,
          switches={9,10},
          switch labels={9/$y_i$,10/\raisebox{-3mm}{$\hat z$}},
        ]{16}
          \shiftedge{8/cfnavy,11/cfnavy,13/cfnavy,16/cfnavy}[$\vec f$]
          \shiftchainbounds{0}{-3.7}{16.2}{0.9}
        \end{shiftchain}
      }
      \only<2->{%
      \begin{shiftchain}[
        column sep=3.5mm, row sep=5mm, dot radius=2.7pt,
        labels=false, label font=\tiny,
        switches={7,9,10,14},
        switch labels={7/$y_{i-1}$,9/$y_i$,10/\raisebox{-3mm}{$\hat z$},14/$y_{i+1}$},
      ]{16}
        \node at (10,0) {$e_3=g_3$};
        \shiftedge{4/cfred,6/cfred,10/cfred,16/cfred}[$\vec g$]
        \shiftedge{8/cfnavy,11/cfnavy,13/cfnavy,16/cfred}[$\vec f$]
        \shiftchainbounds{0}{-3.7}{16.2}{0.9}
      \end{shiftchain}
      }
    \end{column}
    \begin{column}{0.48\textwidth}
      \centering\scriptsize
      \textbf{Gap $(f_2,f_3)$}\par\smallskip
      \only<1>{%
        \begin{shiftchain}[
          column sep=3.5mm, row sep=5mm, dot radius=2.7pt,
          labels=false, label font=\tiny,
          switches={9,10},
          switch labels={9/$y_i$,10/\raisebox{-3mm}{$\hat z$}},
        ]{16}
          \shiftedge{8/cfnavy,9/cfnavy,12/cfnavy,16/cfnavy}[$\vec f$]
          \shiftchainbounds{0}{-3.7}{16.2}{0.9}
        \end{shiftchain}
      }
      \only<2->{%
      \begin{shiftchain}[
        column sep=3.5mm, row sep=5mm, dot radius=2.7pt,
        labels=false, label font=\tiny,
        switches={7,9,10,14},
        switch labels={7/$y_{i-1}$,9/$y_i$,10/\raisebox{-3mm}{$\hat z$},14/$y_{i+1}$},
      ]{16}
        \node at (10,0) {$e_3=g_3$};
        \shiftedge{4/cfred,6/cfred,10/cfred,16/cfred}[$\vec g$]
        \shiftedge{8/cfnavy,9/cfnavy,12/cfnavy,16/cfred}[$\vec f$]
        \shiftchainbounds{0}{-3.7}{16.2}{0.9}
      \end{shiftchain}
      }
    \end{column}
  \end{columns}
\end{frame}

\begin{frame}[t]{Z2: the pair $\hat z,y_{i+1}$ is impossible}
  \footnotesize\setstretch{1}
  \tcbset{before skip=4pt,after skip=4pt}
  \claimreminder{\Yfive\Ythree\Ztwo}
  \begin{lemma}{Claim}
    No new monochromatic edge has only $\hat z,y_{i+1}$.
  \end{lemma}
  \begin{proof}{Proof}
    \begin{itemize}[itemsep=4pt]
      \item<1-> $y_i$ is outside $\vec f$, so
        $e_3^-=y_i<f_1<\hat z=e_3^+$. Thus $f_1=e_3$.
      \item<2-> Both switches lie in the first gap: $y_{i+1}<f_2$.
      \item<3-> Y5 gives $f_1=e_3<x_i$; Y3 gives $f_2<y_{i+1}$: a contradiction.
      \item<4-> Both possible pairs are excluded. No new monochromatic edge appears,
        and $\vec e$ is repaired: Z2.
    \end{itemize}
  \end{proof}
  \begin{center}
    \begin{shiftchain}[
      column sep=6mm, row sep=5mm, dot radius=3pt,
      labels=false, label font=\scriptsize,
      switches={5,6,9/dashed,12},
      switch labels={5/$y_i$,6/\raisebox{-3mm}{$\hat z$},9/$x_i$,12/$y_{i+1}$},
    ]{16}
      \node at (6,0) {$f_1=e_3$};
      \node at (14,0) {$f_2$};
      \shiftedge{6/cfred,14/cfred,15/cfred,16/cfred}[$\vec f$]
      \uncover<3->{\draw[<-,cfred,thick] (10,0.8) -- (14,0.8)
        node[midway,above] {$f_2<y_{i+1}$};}
      \shiftchainbounds{-2}{-2.7}{16.7}{1.5}
    \end{shiftchain}
  \end{center}
\end{frame}

\begin{frame}[t]{Z3: surviving edges start after the new switch}
  \footnotesize\setstretch{1}
  \tcbset{before skip=4pt,after skip=4pt}
  \claimreminder{\Ysix\Ztwo\Zthree\switchparity}
  \begin{lemma}{Claim}\Zthree\end{lemma}
  \begin{proof}{Proof}
    \begin{itemize}[itemsep=4pt]
      \item<1-> \textbf{Base.} Y6 gives the $(2,3)$ pattern; $Z(0)=\varnothing$.
      \item<2-> \textbf{Step.} By Z2 and induction, a surviving $\vec f$ was a $(2,3)$-edge,
        and every old $z$ lies before $f_1$.
      \item<3-> It cannot contain $\hat z$: its two old switches plus $\hat z$ violate Parity.
      \item<4-> Minimality gives $\vec e<\vec f$, hence $\hat z<e_4\leq f_4$.
        Thus $\hat z<f_1$; the two switches of $\vec f$ stay in $(f_2,f_3)$.
    \end{itemize}
  \end{proof}
  \begin{center}
    \uncover<2->{%
      \begin{shiftchain}[
        column sep=6mm, row sep=5mm, dot radius=3pt,
        labels=false, label font=\scriptsize,
        switches={3,5,6,10,12},
        switch labels={3/$y_{i-1}$,5/$y_i$,6/\raisebox{-3mm}{$\hat z$},10/$y_{i+1}$,12/$y_{i+2}$},
      ]{16}
        \shiftedge{1/cfred,2/cfred,6/cfred,9/cfnavy}[$\vec e$]
        \shiftedge{7/cfnavy,9/cfnavy,13/cfnavy,16/cfnavy}[$\vec f$]
      \end{shiftchain}
    }
  \end{center}
\end{frame}

\begin{frame}[t]{Z4: the new switch brackets the third vertex}
  \footnotesize\setstretch{1}
  \tcbset{before skip=4pt,after skip=4pt}
  \claimreminder{\Yfive\Zfour}
  \begin{lemma}{Claim}\Zfour\end{lemma}
  \begin{columns}[T,onlytextwidth]
    \begin{column}{0.61\textwidth}
      \begin{proof}{Proof}
        \begin{itemize}[itemsep=5pt]
          \item<1-> \textbf{Base.} $Z(0)=\varnothing$.
          \item<2-> \textbf{Step.} Every older inserted switch keeps its witness:
            the $Y$-switches stay fixed.
          \item<3-> Y5 gives $y_i=e_3^-<x_i$.
            The next possible switch position is $e_3^+$, so $\hat z=e_3^+\leq x_i$.
          \item<4-> Take $v=e_3$ and $j=i$.
        \end{itemize}
      \end{proof}
    \end{column}
    \begin{column}{0.36\textwidth}
      \centering\scriptsize\setstretch{1}
      {\let\medskip\smallskip \Zexamplefour}
    \end{column}
  \end{columns}
\end{frame}

\begin{frame}[t]{Z1--Z4: termination gives a proper 2-colouring}
  \footnotesize\setstretch{1}
  \tcbset{before skip=4pt,after skip=4pt}
  \claimreminder{\Ztwo\Zthree}
  \begin{lemma}{Claim}$c_Z$ is a proper 2-colouring.\end{lemma}
  \begin{proof}{Proof}
    \begin{itemize}[itemsep=4pt]
      \item<1-> Z2: the finite set of monochromatic edges strictly shrinks at every step.
      \item<2-> At termination no $(2,3)$-edge remains. By Z3, no monochromatic edge remains.
    \end{itemize}
  \end{proof}
  \begin{center}
    \uncover<2->{%
      \begin{shiftchain}[
        column sep=6mm, row sep=3.5mm, dot radius=3pt,
        label font=\scriptsize,
        switches={3,5,6,9},
        switch labels={3/$y_1$,5/$y_2$,6/\raisebox{-3mm}{$z_1$},9/$y_3$},
      ]{10}
        \shiftvertexcolor{1,2,3,6,10}{cfred}
        \shiftvertexcolor{4,5,7,8,9}{cfnavy}
        \shiftedge{1,2,3,4}[$\vec e_1$]
        \shiftedge{2,3,6,7}[$\vec e_2$]
        \shiftedge{4,5,6,7}[$\vec e_3$]
        \shiftedge{4,5,6,10}[$\vec e_4$]
        \shiftedge{7,8,9,10}[$\vec e_5$]
      \end{shiftchain}
    }
  \end{center}
\end{frame}

\begin{frame}[t]{Theorem: every $r\geq4$ shift-chain is 2-colourable}
  \footnotesize\setstretch{1}
  \claimreminder{\Xone\Xthree\Ysix\Ztwo\Zthree}
  \begin{theorem}{Theorem (Smith, 2026)}
    All $r$-shift-chains of uniformity at least $r_0=4$ are properly 2-colourable.
  \end{theorem}
  \begin{itemize}[itemsep=5pt]
    \item<1-> \textbf{X:} one or two switches per edge;
      monochromatic edges are $(3,4)$-edges.
    \item<2-> \textbf{Y:} move switches left;
      only $(2,3)$-edges can stay monochromatic.
    \item<3-> \textbf{Z:} insert switches;
      the monochromatic-edge set strictly shrinks to empty.
    \item<4-> For $r>4$, keep the first four vertices of each edge.
      These form a 4-uniform shift-chain; its proper colouring also colours every original edge properly.
  \end{itemize}
\end{frame}
